%HOMEWORK template for M 325K
\documentclass{article}
\headheight=7pt         \topmargin=14pt
\textheight=574pt       \textwidth=445pt
\oddsidemargin=18pt     \evensidemargin=18pt

%Begin package import

\usepackage{amsmath, amssymb, amsthm}
\usepackage{mathtools}
\usepackage{pdfpages}
\usepackage{tikz-cd}

%End package import
%Begin custom commands

\newcommand{\C}{\mathbb{C}}
\newcommand{\R}{\mathbb{R}}
\newcommand{\Q}{\mathbb{Q}}
\newcommand{\Z}{\mathbb{Z}}
\newcommand{\N}{\mathbb{N}}
\newcommand{\abs}[1]{\lvert #1\rvert}
\newcommand{\RM}[1]{\mathrm{#1}}
\newcommand{\BF}[1]{\textbf{#1}}
\newcommand{\e}{\epsilon}
\newcommand{\ceil}[1]{\lceil{#1}\rceil}
\newcommand{\floor}[1]{\lfloor{#1}\rfloor}
\newcommand{\rarr}[0]{\rightarrow}
\newcommand{\IM}[1]{\text{Im}(#1)}
\newcommand{\Hom}[0]{\text{Hom}}
\newcommand{\Pow}[1]{\text{Pow}(#1)}

%End custom commands
%Begin custom theorems

\newtheorem{innercustomgeneric}{\customgenericname}
\providecommand{\customgenericname}{}
\newcommand{\newcustomtheorem}[2]{%
\newenvironment{#1}[1]
{%
\renewcommand\customgenericname{#2}%
\renewcommand\theinnercustomgeneric{##1}%
\innercustomgeneric%
}
{\endinnercustomgeneric}
}

\newcustomtheorem{THRM}{Theorem}
\newcustomtheorem{LEM}{Lemma}
\newcustomtheorem{EX}{Exercise}
\newcustomtheorem{COR}{Corollary}
\newcustomtheorem{PROB}{Problem}
\newcustomtheorem{claim}{Claim}

\newenvironment{solution}
{\renewcommand\qedsymbol{$\blacksquare$}\begin{proof}[Solution]}
{\end{proof}}

\newenvironment{definition}
{\renewcommand\qedsymbol{$\blacksquare$}\begin{proof}[Definition]}
{\end{proof}}

%End custom theorems
\begin{document}

\title{Jordan Form 1}
\author{Arham Lodha}%put your name here
\date{October 07, 2023}%enter the due date here
\maketitle

\begin{PROB}{1}\label{Problem 1.}
    We say that non-zero v is a "generalized eigenvector" for A if $(A - fI)^k v =0$ for some f in F and some $k > 0$. Show that f is unique, ie if v in the kernel of $(A - fI)^s$ and $(A - gI)^t$, then $f =g$. Hint: Let $S = A - fI, T = A -g I$ explain why some high power of $S-T$ will kill v. Conclude $f = g$.
\end{PROB}
\begin{proof}
    Suppose $v \in \ker((A - fI)^s) \cap \ker((A - gI)^t)$. Let $S = A - fI$ and $T = A - gI$. Let $r = s + t$. $ST = (A - fI)(A - gI) = A^2 - gAI - fIA + fgI^2 = A^2 - (g + f)A + fgI$ and $TS = (A - gI)(A - fI) = A^2 - fIA - gAI + fgI^2 = A^2 - (g + f)A + fgI$. Thus $ST = TS$. Since S and T commute we can use the binomial theorem. Thus $(S - T)^r = \sum_{n=0}^{r}\left({r \choose n} S^{n}T^{r - n}\right)$. For all n such that $r - n \geq t$, $T^{r - n}v = T^{r - n - t}T^{t}v = T^{r - n - t}(0) = 0$. For all n such that $r - n < t$, then $S^nT^{r-n} = T^{r-n}S^n$ (by commutativity above) and $s + t - n < t$ and $s - n < 0$ thus $n > s$ thus $S^{n}v = S^{n - s}S^{s}v = S^{n - s}(0) = 0$ thus $T^{r-n}S^nv = 0$. Thus $\forall n \in [0, \ldots, r]$ $S^nT^{r - n}v = 0$. Thus $(S - T)^rv = \sum_{n=0}^{r}\left({r \choose n} S^{n}T^{r - n}\right)v = \sum_{n=0}^{r}\left({r \choose n} S^{n}T^{r - n}v\right) = 0$. However $S - T = (g - f)I$. Then $((g-f)I)^rv = (g-f)^rIv = (g-f)^rv = 0 = 0v$. Thus $(g-f)^r = 0$ and $g-f = 0$ and $g = f$.
\end{proof}

\bigskip

\begin{PROB}{2a}\label{Problem 2a.}
    In fact, we can take k = dim(V):  Let $S = A - fI$, and $W_i := ker S^i$. Explain why these subspaces "increase" with i, ie we have $0 \subset W_1 \subset W_2 \subset ... $.
\end{PROB}
\begin{proof}
    Base Case: Let $k = 1$. Then by definition of a subspace, $0 \subset W_1$.

    Inductive Hypothesis: Assume that $W_k \subset W_{k+1}$

    Inductive Space: Let $W_{k + 2}:= \ker S^{i + 2}$. $\forall v \in W_{k + 1}$ then $S(S^{k + 1}v) = S^{k+2}v$ but since $v \in \ker S^{k + 1}$ then $S^{k + 1}v = 0$ thus $S^{k + 2}v = 0$. Thus $v \in W_{k + 2}$. Thus $W_{k+1} \subset W_{k+2}$. 

    Thus by induction, $0 \subset W_1 \subset W_2 \subset ... $.
\end{proof}
\begin{PROB}{2b}\label{Problem 2b.}
    Show that if $W_i = W_{i+1}$ then $W_{i+1} = W_{i+2}$, and thus all the $W_k$ for $k \geq i$ are the same. 
\end{PROB}
\begin{proof}

    Suppose $W_i = W_{i + 1}$.
    Base Case: By the previous proof we know $W_{i+1} \subset W_{i+2}$. $\forall v \in W_{i + 2}$, Thus $S^{i + 2}v = S^{i + 1}(Sv)$, thus $Sv \in \ker S^{i+1} = W_{i + 1} = W_{i} = \ker S^i$. Thus $Sv \in \ker S^{i}$. Thus $S^{i}Sv = S^{i + 1}v = 0$ and $v \in W_{i + 1}$. Thus $\forall v \in W_{i + 2}$ then $v \in W_{i + 1}$. Thus $W_{i + 1} \subset W_{i + 2}$ and $W_{i + 1} = W_{i + 2}$. 

    Inductive Hypothesis: For a positive integer $k > 0$, assume $W_{i + k} = W_{i + k + 1}$.

    Inductive Step: By the previous proof we know $W_{i + k +1} \subset W_{i +k+2}$. $\forall v \in W_{i + k + 2}$. By defintion of kernel, $S^{i + k + 2}v = 0 = S^{i + k + 1}Sv = 0$. Thus $Sv \in W_{i + k + 1}$. By the inductive hypothesis, since $W_{i + k + 1} = W_{i + k}$. Thus $Sv \in W_{i + k}$ thus $S^{i+k}Sv = S^{i+k+1}v = 0$. Thus $v \in W_{i + k + 1} = W_{i + k}$. Thus $W_{i + k + 2} = W_{i + k + 1} = W_{i + k}$.

    Thus by induction, $\forall k \in \Z$ s.t $k > 0$ then $W_i = W_{i + k}$.
\end{proof}

\begin{PROB}{2c}\label{Problem 2c.}
    Now explain why $W_k = W_{\dim V}$ for sufficiently big k, and conclude the generalized eigenvectors are the (non-zero) vectors in $\ker (A - fI)^{\dim V}$. Now we define the generalized eigenspace, for f, to be $GV_f(A) := \ker (A - fI)^{\dim V}$. 
\end{PROB}
\begin{proof}
    Since $W_i \subseteq V$ are subspaces of V, $\forall i \in \N$. The $\dim W_i \leq \dim V$. Thus $\dim W_{\dim V} \leq \dim V$. If $\dim W_{\dim V} < \dim V$, then there is a $0 \geq k < \dim V$ s.t $W_{k} = W_{k + 1}$ because otherwise $\dim W_{k + 1} > \dim W_{k}$ for all k. and then $\dim W_{\dim V} = \dim V$ forming a contradiction with the assumption. Thus $\exists k$ s.t $W_k = W_{k+1}$. By previous proof, $W_k = W_{\dim V}$ because $\dim V > k$. However if $\dim W_{\dim V} = \dim V$ and since $W_{\dim V} \subseteq W_{\dim V+1}$ thus $\dim W_{\dim V} \leq \dim W_{\dim V + 1} \leq \dim V$ thus $\dim W_{\dim V + 1} = \dim V$ thus $\dim W_{\dim V} = \dim W_{\dim V + 1}$ and $W_{\dim V} = W_{\dim V + 1}$. Thus $k = \dim V + 1$. By the previous theorem, $\forall n \geq k$ $W_{\dim V} = W_{n}$.

    Thus $W_{\dim V}$ contains all $W_i$ for $i \leq \dim W$ and is equal to all $W_i$ for $i \geq \dim W$. Thus all the generalized eigenvectors are the nonzero vectors in $W_{\dim V}$. 
\end{proof}

\begin{PROB}{2d}\label{Problem 2d.}
    Show this is non-zero iff f is an eigenvalue. 
\end{PROB}
\begin{proof}
    $(\implies)$: $GV_f(A) \neq 0$. 

    \begin{enumerate}{}
        \item If there exists no $k \in [1, \ldots \dim V]$ such that $\exists v \in GV_f(A)$ s.t $v \not \in W_k$ then $W_1 = GV_f(A)$. Thus since $GV_f(A) \neq 0$, then $W_1 \neq 0$ $W_1$ is by definition a eigenspace of A. Thus $f$ must be a eigenvalue.
        \item If there exists a $k \in [1, \ldots \dim V]$ such that $\exists v \in GV_f(A)$ s.t $v \not \in W_k$. Thus $\forall v \in W_k$, $v \not \in V_{k - 1}$ because if otherwise $k = k- 1$ which is a contradiction. Thus $\forall v \in W_k$, $(A - fI)^kv = 0$ and $(A - fI)((A - fI)^{k-1}v) = 0$. Since $(A - fI)^{k-1}v \neq 0$, then $f$ must be a eigenvalue by definition.
    \end{enumerate}
    
    $(\impliedby)$: Suppose $f$ is a eigenvalue. Thus there is a nonzero eigenvector $v$ such that $(A - fI)v = 0$ by definition of eigenvalues and eigenvectors. Thus $v \in W_1$. Since $W_1 \subset W_{\dim V}$ by 2a. $v \in W_{\dim V} = GV_f(A)$. Thus $GV_f(A) \neq 0$. 
\end{proof}

\bigskip

\begin{PROB}{3}\label{Problem 3.}
    What do you expect is true about the generalized Eigenspaces? Lets say $GV_f(A)$ is the kernel of $(A - fI)^{dim V}$. What do you think is true of the map  $\oplus_{f \in F} GV_f(A) \to V$ (the addition map). Prove it -- the argument is similar to the eigenspace case, just a bit more involved.
\end{PROB}
\begin{solution}
    Since $GV_f(A) \subset V$ is a subspace and thus $\dim(GV_f(A)) \leq \dim V$.
\end{solution}
\begin{claim}{3}
    $\oplus_{f \in F} GV_f(A) \to V$ is injective. 
\end{claim}
\begin{proof}
    \textbf{Base Case (k = 2): } For $f, g \in F$ where $f \neq g$, $\forall v \in GV_{f}(A)$ and $\forall w \in GV_{g}(A)$. Let $v + w = 0$. $A^{\dim(V)}(v + w) = 0$ and $A^{\dim(V)}v + A^{\dim(V)}w = 0$. Since $v \in GL_{f}(A)$ and $w \in GL_{g}(A)$. Thus $fv + gw = 0$. $f(v + w) = 0$. Thus $fv + fw = 0$. Thus $fv + gw - fv - fw = (g - f)w = 0$. But since $f \neq g$, then $g - f \neq 0$ thus $w = 0$. But since $v + w = 0$, then $v + 0 = 0$ and $v = 0$.

    \textbf{Inductive Hypothesis: } If $\sum_{i = 1}^k v_i = 0$ for $v_i \in GV_{f_i}(A)$ with $f_i$ distinct implies $v_i = 0$.

    \textbf{Inductive Step: } Suppose $\sum_{i = 1}^{k + 1} v_i = 0$ for $v_i \in GV_{f_i}(A)$ with $f_i$ distinct. Then $\sum_{i = 1}^{k + 1} v_i = \sum_{i = 1}^k v_i + v_{k+1} = 0$. By our inductive hypothesis, $0 + v_{k + 1} = v_{k + 1} = 0$. 

    Thus by induction $\sum_{i = 1}^{k + 1} v_i = 0 \implies v_i = 0$. Thus the kernel of $\oplus_{f \in F}$ is only the 0 vector. Thus the $\oplus_{f \in F}$ is injective. Thus $\oplus_{f \in F} GV_f(A) \subset V$ and $\dim(\oplus_{f \in F} GV_f(A)) \leq \dim V$
\end{proof}

\bigskip

\begin{PROB}{4}\label{Problem 4.}
    Now lets prove that this map is always an isomorphism, ie V is always the direct sum of the generalized eigenspaces (assuming that the characteristic polynomial has all its roots in F). To do this use the Cayley-Hamilton theorem (which we will prove next term, if not before then): namely that any matrix satisfies its characteristic polynomial, together with TKL (that kids lemma, about the dim of the kernel of a product).
\end{PROB}
\begin{proof}
    The generalized eigenvalues of $A$ is $f_1, \ldots, f_n$. 
    
    The characteristic polynomial where $m = \dim V$, $p_A(x) = (x - f_1)^{m}\ldots(x - f_n)^{m}$. By Cayley Hamilton Theorem, $p_A(A) = 0$. Thus $\ker p_A(A) = V$ thus $\dim \ker p_A(A) = \dim V$. By TKL, $\dim V = \dim \ker p_A(A) \leq \dim \ker (x - f_1)^{m} + \ldots + \dim \ker (x - f_n)^{m} = \dim GV_{f_1}(A) + \ldots + \dim GV_{f_n}(A)$. Thus $\dim V \leq \dim GV_{f_1}(A) + \ldots + \dim GV_{f_n}(A)$. But $\oplus_{f \in F} GV_f(A)$ is a subspace of $V$, thus $\dim V \geq \dim GV_{f_1}(A) + \ldots + \dim GV_{f_n}(A)$. Thus $\dim V = \dim GV_{f_1}(A) + \ldots + \dim GV_{f_n}(A)$. Thus oplus is a bijection and thus a isomorphism.
    
\end{proof}


\end{document}